\chapter{Implementation}
\label{ch:impl}

This chapter describes how the design of \Cref{ch:design} is realised
in code: the libraries chosen at each layer, the repository structure,
the concrete configuration values, and the engineering choices that
keep the system reproducible within the compute budget of a single
consumer GPU.
The conceptual architecture in \Cref{ch:design} is intentionally
library-agnostic; this chapter fills in the specific versions, magic
numbers, and trade-offs that follow from the hardware and the toolchain.
Computational-complexity arguments and speed optimisations are
collected in \Cref{sec:impl-performance}.

The realised system, at the scale reported in \Cref{ch:results},
renders \num{87677} PrIMuS lines, adds \num{26168} header-less twin
samples (\code{\_\_nh}), and trains on \num{71297} in-domain pairs
after the filters of \Cref{sec:design-data} (from \num{113845}
augmented images including twins), with a vocabulary of \num{80}
LMX symbols (\num{77} content tokens in
\code{data/vocab/primus\_lmx.txt} plus blank, pad, and unk), and
converges in
\num{37}~epochs of approximately \SI{8}{\minute} each on the target
hardware (measured as the mean of epochs 2--49 from the training log;
the first epoch is longer due to data-loader warm-up).
The final checkpoint occupies approximately \SI{165}{\mega\byte}.

% ─────────────────────────────────────────────────────────────────────────────
\section{Technology Stack}
\label{sec:impl-stack}

\Cref{tab:tech-stack} lists the principal libraries used in each layer
of the system.
The selection criterion in every case was the smallest, most
widely-used library compatible with the design constraints stated in
\Cref{ch:design}; no component was chosen for novelty.


\begin{table}[htbp]
\centering
\small
\setlength{\tabcolsep}{5pt}
\renewcommand{\arraystretch}{1.35}
\begin{tabularx}{\textwidth}{@{} l l L @{}}
\toprule
\textbf{Layer} & \textbf{Library / version} & \textbf{Role} \\
\midrule
Deep learning            & PyTorch $\geq 2.10$~\cite{paszkeImperativeStyleHighPerformance2019} & Model, loss, optimiser, mixed-precision training. \\ \tfgrowsep
CNN backbone             & \code{torchvision} $\geq 0.25$~\cite{torchvision2016} & ResNet18 module definition. \\ \tfgrowsep
Image processing         & OpenCV $\geq 4.13$~\cite{bradskiOpenCVLibrary2000} & Binarisation, deskewing, morphological staff detection. \\ \tfgrowsep
PDF rasterisation        & PyMuPDF (\code{fitz}) $\geq 1.27$~\cite{artifexPyMuPDF} & Render PDF pages to grayscale arrays at configurable DPI. \\ \tfgrowsep
Music rendering          & LilyPond $\geq 2.25$ + LilyJAZZ~\cite{nienhuysLilyPondSystem2003,hammerleOpenLilyPondFontsLilyJAZZ} & Generate training PNGs from PrIMuS semantic files. \\ \tfgrowsep
Data augmentation        & \code{albumentations} $\geq 2.0$\,\cite{AlbumentationsFastFlexible} & Scan-simulation transform pipeline. \\ \tfgrowsep
Numerics                 & NumPy $\geq 2.4$~\cite{harrisArrayProgramming2020} & Tensor / array interop. \\ \tfgrowsep
Text-region OCR          & EasyOCR $\geq 1.7$~\cite{jaidedaiEasyOCR} & Pre-CRNN staff-reject gate (text-density signal). \\ \tfgrowsep
Web service              & FastAPI $\geq 0.115$ + Uvicorn $\geq 0.32$~\cite{ramirezFastAPI} & REST API wrapping the inference pipeline. \\ \tfgrowsep
Package manager          & Poetry $\geq 2.0$ & Dependency isolation, deterministic environment. \\
\bottomrule
\end{tabularx}
\caption{Principal libraries used by the system and the version pins
recorded in \code{pyproject.toml}.  Poetry's \code{poetry.lock} fixes
every transitive dependency, so the environment is reproducible
across machines with a single \code{poetry install}.}
\label{tab:tech-stack}
\end{table}

The runtime target is a single NVIDIA RTX~3060 with \SI{12}{\giga\byte}
of VRAM, which is the hardware available to the author for the
duration of the thesis.
All architectural choices in \Cref{ch:design} were taken with this
budget in mind; the implementation does not require multi-GPU support,
distributed training, or specialised cloud infrastructure.
The Python version is pinned at \num{3.14}; runtime
dependencies are listed in \code{pyproject.toml} and locked in
\code{poetry.lock}.

% ─────────────────────────────────────────────────────────────────────────────
\section{Repository Structure}
\label{sec:impl-repo}

The project source tree is organised by responsibility rather than by
file type, with three structural rules.
First, all music-notation mapping tables (clef IDs, key-signature
encodings, duration tokens) live in a single module,
\code{CRNN\_CTC/lilypond\_render.py}, imported by every module that
needs them; this stops the rendering pipeline and the label converter
drifting out of sync.
Second, the inference pipeline (\code{omr\_pipeline/}) depends on the
training code only through the \code{Vocabulary} and \code{CRNN}
classes, so it can be deployed without the training tree.
Third, \code{cli.py} is the only entry point: every multi-step
operation (data generation, training, evaluation, vocab building,
threshold calibration) is a subcommand sharing the same flag names and
defaults.

\begin{quote}\small
\begin{verbatim}
src/
  cli.py                      unified CLI (12 subcommands)
  CRNN_CTC/                   model, training, vocab, evaluation
    model.py                  CRNN-CTC architecture
    vocab.py                  Vocabulary (blank=0, pad=1, unk=2)
    dataset.py, config.py     PrIMuS dataset + Config dataclass
    train.py, evaluate.py     training loop and SER evaluation
    lilypond_render.py        single source of truth for clef/
                              key/duration → LilyPond tables
    training_utils.py         shared train / eval helpers
    chord_dataset.py          synthetic chord-strip dataset
    chord_train.py            chord CRNN training from scratch
    chord_finetune.py         fine-tune chord CRNN on real strips
  data_processing/            dataset generation and augmentation
    generate_realbook.py      PrIMuS -> LilyJAZZ render
    semantic_to_lmx.py        PrIMuS .semantic -> LMX label
    augment_scanned.py        scan-simulation augmentation
    chord_render.py           synthetic chord-strip renderer
    generate_chord_crops.py   bulk chord-strip generation
    extract_real_chord_strips.py    pre-label strips from PDFs
    extract_real_music_strips.py    pre-label staff strips
  omr_pipeline/               inference pipeline
    pipeline.py               run_pipeline() entry point
    preprocess.py             load / binarise / deskew
    staff_detect.py           morphology -> 5-line clusters
    staff_reject.py           hybrid pre/post-CRNN rejection gate
    inference.py              CRNN-CTC music recognition
    chord_recognizer.py       chord OCR
    chord_postprocess.py      jazz-chord grammar filter
    grammar_fix.py            LMX token sequence validator
  api/                        FastAPI web server
\end{verbatim}
\end{quote}

% ─────────────────────────────────────────────────────────────────────────────
\section{Data Pipeline Implementation}
\label{sec:impl-data}

The data pipeline of \Cref{sec:design-data} is implemented as five
offline stages orchestrated by \code{pipeline} and \code{pipeline-train}:
\code{render}, \code{convert}, header-less twin generation
(\code{generate\_headerless\_twins.py}), \code{augment}, and \code{vocab}.
Incremental runs skip stages whose outputs are already up to date;
\code{--force-all} rebuilds everything from scratch.
\Cref{fig:data-pipeline-impl} shows the full flow.
Every stage uses the Python \code{multiprocessing} module for CPU
parallelism, with the worker-pool size set by a CLI flag (default 8 on
the target machine).
The choice of process-based parallelism over threading follows from the
Global Interpreter Lock; the trade-off is discussed in
\Cref{sec:impl-performance}.

\begin{figure}[ht]
\centering
\begin{tikzpicture}[
  >={Stealth[round]},
  io/.style={
    draw, shape=regular polygon, regular polygon sides=6,
    shape border rotate=90, inner sep=2pt, align=center,
    font=\scriptsize\sffamily
  },
  stage/.style={
    draw, rectangle, rounded corners=3pt,
    minimum width=4.8cm, minimum height=0.95cm,
    align=center, font=\small\sffamily,
    inner sep=4pt
  },
  flow/.style={->, thick},
  note/.style={font=\scriptsize\itshape, text=black!55},
  node distance=0.6cm
]
  % ── Nodes ────────────────────────────────────────────────────────────────
  \node[io]                       (raw)    {PrIMuS\\raw data};
  \node[stage, below=of raw]      (render) {LilyJAZZ render\\\footnotesize\ttfamily generate\_realbook.py};
  \node[stage, below=of render]   (conv)   {Semantic-to-LMX conversion\\\footnotesize\ttfamily semantic\_to\_lmx.py};
  \node[stage, below=of conv]     (twins)  {Header-less twin generation\\\footnotesize\ttfamily generate\_headerless\_twins.py};
  \node[stage, below=of twins]    (aug)    {Scan-simulation augmentation\\\footnotesize\ttfamily augment\_scanned.py};
  \node[stage, below=of aug]      (vocab)  {Vocabulary construction\\\footnotesize\ttfamily cli.py~vocab};
  \node[io, below=of vocab]       (split)  {Train / val\\/ test splits};

  % ── Output annotations ───────────────────────────────────────────────────
  \node[note, right=0.9cm of conv]  (n1) {PNG\,+\,\texttt{.lmx} label pairs};
  \node[note, right=0.9cm of twins] (n2) {\texttt{\_\_nh} header-less variants};
  \node[note, right=0.9cm of aug]   (n3) {distorted copies in \texttt{scanned/}};
  \node[note, right=0.9cm of vocab] (n4) {\texttt{data/vocab/primus\_lmx.txt}};

  % ── Arrows ───────────────────────────────────────────────────────────────
  \draw[flow] (raw)    -- (render);
  \draw[flow] (render) -- (conv);
  \draw[flow] (conv)   -- (twins);
  \draw[flow] (twins)  -- (aug);
  \draw[flow] (aug)    -- (vocab);
  \draw[flow] (vocab)  -- (split);

  % ── Outer system box + bold label ────────────────────────────────────────
  % fit= is unreliable when notes extend rightward; use explicit coordinates.
  % twins.west gives the leftmost stage edge; n4.east the rightmost note edge.
  \begin{scope}[on background layer]
    \draw ([xshift=-0.4cm, yshift=0.4cm]  raw.north   -| twins.west)
          rectangle
          ([xshift=0.5cm,  yshift=-0.4cm] split.south -| n4.east);
  \end{scope}
  % title north must be below the rectangle bottom (split.south - 0.4cm);
  % use -0.45cm so there is a small gap, matching the other pipeline figures
  \node[anchor=north, font=\small\bfseries\sffamily, inner sep=3pt]
    at ($([xshift=-0.4cm] split.south -| twins.west)!0.5!([xshift=0.5cm] split.south -| n4.east)+(0,-0.45cm)$)
    {Data preparation pipeline};

\end{tikzpicture}
\caption[Data preparation pipeline.]{Five offline stages that convert PrIMuS
  raw data into training-ready image/label pairs.
  The \code{pipeline} and \code{pipeline-train} subcommands run all five
  stages in sequence; incremental runs skip any stage whose outputs are already
  up to date.}
\label{fig:data-pipeline-impl}
\end{figure}

\subsection{Render: PrIMuS to LilyJAZZ}

The \code{render} subcommand walks the PrIMuS package directories and,
for each \code{.semantic} file, normalises the clef to G2 (transposing
pitches to the treble-clef equivalent when the source clef is C1, C2,
or F3), constructs a LilyPond source file, and invokes the
\code{lilypond} subprocess at \SI{200}{DPI}.
The output PNG is auto-cropped to remove white margins.

Four per-sample parameters are drawn independently to cover the visual
range described in \Cref{sec:design-data}.
The global staff size is sampled uniformly from
$\{17, 18, 19, 20, 21, 22\}$ points, which varies notehead size and
horizontal note spacing.
Common-time and cut-time signatures are rendered as numerals with
probability \num{0.5} and as the corresponding C glyph otherwise.
Cautionary accidentals are added with probability \num{0.5}.
Beaming is beat-grouped in \SI{85}{\percent} of samples and left
unbeamed in the remaining \SI{15}{\percent}, matching the typical ratio
in Real Book editions.

The LilyJAZZ font stylesheet is not bundled with LilyPond and must be
installed manually into the LilyPond \code{share} directory before
rendering can begin; the installation procedure is in
\code{docs/data\_pipeline.md}.
LilyPond~$\geq$~2.25 also requires a small patch to LilyJAZZ: the
deprecated \code{set-global-fonts} call is rewritten using the
\code{property-defaults.fonts.*} API introduced in that version.

\subsection{Convert: PrIMuS \texttt{.semantic} to LMX}

The \code{convert} subcommand parses each \code{.semantic} file
token by token, applies the display-accidental rules of
\Cref{sec:lmx-format} from the pitch spelling and the active key
signature, and emits an LMX label file co-located with the rendered
PNG.
The conversion is direct: no intermediate MusicXML representation is
produced.

The non-obvious implementation detail is the \term{implicit-key fix}:
PrIMuS omits the \code{keySignature-} token entirely for C-major
samples (no accidentals = default), and without intervention the model
would have to learn to predict \code{key:fifths:0} from a blank key
area without explicit supervision.
The converter detects this case and injects the C-major key token
before the time signature so every training sample carries an
explicit key label.
After the fix, C major accounts for roughly \SI{22.6}{\percent} of the
post-filter corpus (\num{19778} of \num{87677} samples), confirming
that the missing label was a labelling artefact rather than a
distribution effect.

\subsection{Augment: scan simulation}

The \code{augment} subcommand applies the
\code{albumentations}~\cite{AlbumentationsFastFlexible} transform
pipeline of \Cref{tab:aug-pipeline} to every clean image and writes the
result to a parallel directory.
After the library pipeline, six non-library operations are applied in
fixed order: morphological ink dilation (0, 1, or 2 erosion iterations
chosen stochastically), staff-line thickness variation ($p=0.40$),
paper-tone remapping from pure white/black to the paper-like range
$[28, 245]$, a radial vignette ($p=0.55$), a directional
uneven-illumination gradient simulating book-spine shadow ($p=0.45$),
and faint horizontal halftone banding ($p=0.25$).

\begin{table}[htbp]
\centering
\small
\setlength{\tabcolsep}{5pt}
\renewcommand{\arraystretch}{1.35}
\begin{tabularx}{\textwidth}{@{} l c L @{}}
\toprule
\textbf{Transform} & \textbf{Probability} & \textbf{Range / purpose} \\
\midrule
\code{ElasticTransform}         & 0.82 & page warping, staff curvature \\ \tfgrowsep
\code{GridDistortion}           & 0.72 & grid bending \\ \tfgrowsep
\code{OpticalDistortion}        & 0.38 & lens distortion \\ \tfgrowsep
\code{Affine}                   & 0.85 & rotation $\pm 4^\circ$, shear $\pm 1.2^\circ$, scale $0.97$--$1.03$ \\ \tfgrowsep
\code{GaussianBlur}             & 0.62 & $\sigma \in [0.2, 1.1]$ \\ \tfgrowsep
\code{Sharpen}                  & 0.48 & over-sharpened-scanner simulation \\ \tfgrowsep
\code{RandomToneCurve}          & 0.72 & exposure variation \\ \tfgrowsep
\code{GaussNoise}               & 0.78 & sensor noise \\ \tfgrowsep
\code{RandomBrightnessContrast} & 0.70 & brightness $\pm 0.06$, contrast $0.03$--$0.18$ \\ \tfgrowsep
\code{ImageCompression}         & 0.35 & JPEG quality $72$--$92$ \\
\bottomrule
\end{tabularx}
\caption{The \code{albumentations} transform pipeline applied during
the \code{augment} stage.  Probabilities and parameter ranges are
calibrated against a sample of real Real Book scans.}
\label{tab:aug-pipeline}
\end{table}

Representative outputs of the augmentation pipeline are shown in
\Cref{fig:aug-gallery}; the present section focuses on the
implementation details that produce them.

A second, lighter augmentation layer is applied online during training
inside \code{Dataset.\_\_getitem\_\_}: per-sample brightness and
contrast jitter, low-amplitude noise, and a small horizontal
translation, each with probability \num{0.5}.
Without this online jitter, every epoch sees identical pixel grids and
the model overfits the exact augmentations baked into
\code{scanned/}.
The online cost is approximately \SI{50}{\micro\second} per sample,
small enough to be absorbed by the data-loader's worker processes.

\subsection{Vocab: build the token index}

The \code{vocab} subcommand scans every \code{.lmx} file in the
training corpus, collects the unique tokens, sorts them
alphabetically, and writes them to a plain-text file with one token
per line.
The three special tokens (\code{<blank>}, \code{<pad>}, \code{<unk>})
are absent from the file; the \code{Vocabulary} class injects them at
indices 0, 1, and 2 at load time, so line $N$ in the file maps to
index $N+3$.
The rationale for the index layout and the guaranteed pitch/octave
reservations is described in \Cref{sec:lmx-format}.

% ─────────────────────────────────────────────────────────────────────────────
\section{Model Implementation}
\label{sec:impl-model}

\subsection{Module structure}

The model is implemented as three PyTorch \code{nn.Module} subclasses:
a CNN backbone, a recurrent stack, and a thin linear projection head.
The top-level \code{CRNN} class composes them and exposes a single
\code{forward(images, image\_widths)} method that returns log-softmax
probabilities ready for the CTC loss.

\subsection{ResNet18 backbone modifications}

The model uses \code{torchvision.models.resnet18} without ImageNet
pre-trained weights (\code{resnet18(weights=None)}); the rationale
for training from scratch on the single-channel monochrome input is
in \Cref{sec:design-model}.
Two modifications turn the standard ResNet18 into a CRNN backbone.
First, \code{conv1} is replaced with a $7\times7$ convolution accepting
single-channel input.
Second, the strides of three convolutional stages are made asymmetric
so that the height dimension is compressed eight times more
aggressively than the width:

\begin{table}[ht]
\centering
\small
\setlength{\tabcolsep}{8pt}
\renewcommand{\arraystretch}{1.20}
\begin{tabularx}{0.85\textwidth}{@{} l c c L @{}}
\toprule
\textbf{Stage} & \textbf{Stride (h, w)} & \textbf{Output (h, w)} & \textbf{Effect} \\
\midrule
\code{conv1}                  & (2, 1) & $H/2,\ W$    & halve height, keep width \\
\code{maxpool}                & (2, 2) & $H/4,\ W/2$  & halve both \\
\code{layer1}                 & (1, 1) & $H/4,\ W/2$  & residual block \\
\code{layer2}                 & (2, 2) & $H/8,\ W/4$  & halve both \\
\code{layer3}                 & (2, 1) & $H/16,\ W/4$ & halve height \\
\code{layer4}                 & (2, 1) & $H/32,\ W/4$ & halve height \\
\code{AdaptiveAvgPool2d}      & --     & $1,\ W/4$    & collapse height to 1 \\
\bottomrule
\end{tabularx}
\caption{Stride table for the modified ResNet18 backbone.  Cumulatively
the input shape $(B, 1, H{=}128, W)$ becomes
$(B, 512, 1, W/4)$ at the backbone output: the height is collapsed to
one feature vector per column, and the time axis (width) is preserved
with a $4\times$ compression.}
\label{tab:resnet-strides}
\end{table}

After the backbone, the height dimension is squeezed so the BiLSTM
receives a tensor of shape $(W/4, B, 512)$.
The $4\times$ width compression follows the design argument of
\Cref{sec:design-model}: it is the smallest power-of-two ratio at
which every in-scope sample comfortably satisfies the CTC
input-length~$\geq$~label-length constraint.

\subsection{Bidirectional LSTM and projection head}

The recurrent stack is a two-layer bidirectional \code{nn.LSTM} with
\num{256} hidden units per direction (\num{512} total), inter-layer
dropout of \num{0.3}, and the standard PyTorch packed-sequence handling
for variable-width batches.
The output is projected to the vocabulary size by an \code{nn.Linear}
layer, followed by \code{log\_softmax} along the vocabulary dimension;
this is the input shape required by \code{nn.CTCLoss}.

\subsection{CTC loss configuration}

The loss is \code{nn.CTCLoss(blank=0, zero\_infinity=True)}.
The \code{blank=0} convention is a hard contract of the project: the
\code{Vocabulary} class hardcodes \code{<blank>=0}, and every decode
path assumes the blank lives at index $0$.
\code{zero\_infinity=True} silently zeroes the loss contribution of
samples whose label length exceeds the effective input length
($W/4$), which would otherwise yield an infinite loss and a NaN
gradient.
On the in-scope corpus this masking is essentially never triggered
because the multi-staff filter of \Cref{sec:design-data} removes the
samples most at risk; the flag is a safety net rather than a routine
code path.

\subsection{Training loop}

The training loop is a standard single-GPU PyTorch loop.
Each epoch runs the following steps in order:

\begin{enumerate}
    \item Iterate over the training \code{DataLoader} configured with
          \code{num\_workers=10}, \code{pin\_memory=True}, and
          \code{persistent\_workers=True} (the last two avoid the
          re-spawn cost of worker processes between epochs).  Each
          batch is a dictionary of right-padded image tensor, packed
          label tensor, per-sample label lengths, and per-sample image
          widths.
    \item Forward pass through the CRNN inside an
          \code{torch.amp.autocast} context for mixed-precision
          training; the CTC loss itself is computed in single
          precision, the model body in \code{float16}.
    \item Backward pass via a \code{torch.amp.GradScaler} on the
          scaled loss; gradient norm clipped at \num{5.0} to stabilise
          the LSTM.
    \item Optimiser step (\code{AdamW}~\cite{loshchilovDecoupledWeightDecay2019},
          peak learning rate $1{\times}10^{-3}$,
          weight decay $1{\times}10^{-4}$).  The optimiser is
          constructed with \code{fused=True} on CUDA hardware to
          collapse the parameter-update kernel.
    \item Scheduler step (\code{OneCycleLR}~\cite{smithSuperConvergenceVeryFast2019},
          single-cycle cosine annealing, warm-up fraction \num{0.08},
          stepping once per minibatch).
    \item Validation pass at end of epoch: greedy CTC decode on the
          full validation set, SER computed against ground truth, and
          \code{best\_model.pt} written when the validation SER
          improves.
    \item Append the epoch's metrics to \code{training\_log.csv} in
          the model directory.
    \item Stop if SER has not improved for \num{12} consecutive epochs
          (early stopping).
\end{enumerate}

\subsection{Hyperparameter summary}

\Cref{tab:hyperparams} lists every hyperparameter used by the default
training configuration.  All defaults are encoded in the \code{Config}
dataclass and serialised to \code{config.json} at training start, and
they are written into every checkpoint so that evaluation and
inference always use the exact settings that produced the weights.
Because the \code{Config} dataclass is part of the checkpoint contract,
adding, renaming, or removing a field is a backward-incompatible
change for existing checkpoints; this is recorded as
Hard~Constraint~3 of the project's contributor guide.

\begin{table}[htbp]
\centering
\small
\setlength{\tabcolsep}{5pt}
\renewcommand{\arraystretch}{1.25}
\begin{tabularx}{\textwidth}{@{} l l L @{}}
\toprule
\textbf{Hyperparameter} & \textbf{Default} & \textbf{Notes} \\
\midrule
\code{img\_height}               & 128                & fixed; image is resized to this height \\ \tfgrowsep
\code{max\_image\_width}         & 2048               & wider images are clamped \\ \tfgrowsep
\code{max\_source\_height}       & 180                & multi-staff filter threshold (on the pre-resize PNG) \\ \tfgrowsep
\code{backbone}                  & \code{resnet18}    & VGG retained for legacy checkpoints only \\ \tfgrowsep
\code{rnn\_hidden}               & 256                & per direction \\ \tfgrowsep
\code{rnn\_layers}               & 2                  & BiLSTM depth \\ \tfgrowsep
\code{dropout}                   & 0.3                & between LSTM layers \\ \tfgrowsep
\code{batch\_size}               & 16                 & peak VRAM $\approx$ \SI{8.5}{\giga\byte} on RTX~3060 with AMP \\ \tfgrowsep
\code{lr}                        & $1{\times}10^{-3}$ & OneCycleLR peak \\ \tfgrowsep
\code{weight\_decay}             & $1{\times}10^{-4}$ & AdamW \\ \tfgrowsep
\code{warmup\_frac}              & 0.08               & fraction of total steps for warm-up \\ \tfgrowsep
\code{max\_grad\_norm}           & 5.0                & gradient clipping \\ \tfgrowsep
\code{epochs}                    & 60                 & maximum, with early stopping \\ \tfgrowsep
\code{early\_stopping\_patience} & 12                 & epochs without val SER improvement \\ \tfgrowsep
\code{val\_frac}                 & 0.10               & validation split fraction \\ \tfgrowsep
\code{test\_frac}                & 0.10               & test split fraction \\ \tfgrowsep
\code{strip\_header\_prob}       & 0.0                & deprecated/inert; header-less continuation staves are first-class \code{\_\_nh} twin samples (see \Cref{sec:design-training}) \\ \tfgrowsep
\code{online\_aug\_prob}         & 0.50               & in-loader jitter rate (training only) \\ \tfgrowsep
\code{rare\_lmx\_oversample}     & 2                  & multiplier for samples with tie tokens \\ \tfgrowsep
\code{num\_workers}              & 10                 & DataLoader worker processes \\
\bottomrule
\end{tabularx}
\caption{Default hyperparameters of the \code{Config} dataclass.  Every
value is overridable via a CLI flag; the entire \code{Config} is
serialised alongside the checkpoint so that evaluation always uses
the same settings as training.}
\label{tab:hyperparams}
\end{table}

% ─────────────────────────────────────────────────────────────────────────────
\section{Chord OCR Implementation}
\label{sec:impl-chord}

The chord-recognition network reuses the CRNN--CTC machinery of the
melody network but is trained on a different data stream with a
different vocabulary.
The two-stream split was justified in \Cref{sec:design-twostream}; the
remainder of this section reports the engineering details specific to
the chord branch.

\subsection{Synthetic chord-strip generation}

The \code{chord\_render} module is the single source of truth for
chord-notation conventions on the training side: it defines the
canonical set of qualities (major, minor, dominant 7, major 7,
half-diminished, diminished, augmented, suspended) together with their
sampling weights, which are derived from the empirical frequency of
each quality in the Real Book.
The \code{generate\_chord\_crops} script renders strips of one to five
chords each, in the LilyJAZZ font, at the same DPI as the staff
renders.  It then applies an \code{albumentations} pipeline tuned for
the visual conditions of the chord strip (smaller scale, less
geometric distortion, more colour-channel jitter) and an additional
\textit{Real Book clutter} overlay
(\code{\_add\_realbook\_clutter}) that adds artefacts such as
binder-hole shadows, staff-line bleed from the staff below, and
occasional slash-repeat marks.

The default configuration produces \num{12000} training strips and
\num{1200} validation strips; the chord training script consumes them
via a CSV manifest.

\subsection{Chord network architecture and training}

The chord network is the same ResNet18 + BiLSTM + linear-projection
stack as the melody network, but with a \num{29}-class output
(\num{26} character tokens plus blank, pad, unk).
The input image is resized to \SI{64}{px} of height (half the staff
height) before being fed to the network; the same height is used at
inference time.

The network is scaled down relative to the melody CRNN because chord
text is visually simpler than staff notation: the BiLSTM uses
\num{192} hidden units per direction (versus \num{256} for the melody
network), and the CNN applies an additional inter-block dropout of
\num{0.15}.

Training runs in two phases.
The synthetic pretraining phase uses AdamW with a peak learning rate
of $1{\times}10^{-3}$ and OneCycleLR with a \SI{5}{\percent}
warm-up fraction, a batch size of \num{32}, and early stopping with a
patience of \num{8} epochs.
In practice the model reaches a stable Character Error Rate on the
validation strips within ten to fifteen epochs.
The fine-tuning phase starts from the synthetic checkpoint and runs
for up to \num{20} epochs at a reduced peak learning rate of
$2{\times}10^{-4}$; the lower rate prevents the small real-strip set
from overwriting chord representations built during pretraining.
\Cref{tab:chord-hyperparams} lists the full parameter set for both
phases.

\begin{table}[htbp]
\centering
\small
\setlength{\tabcolsep}{5pt}
\renewcommand{\arraystretch}{1.25}
\begin{tabularx}{\textwidth}{@{} l l l L @{}}
\toprule
\textbf{Parameter} & \textbf{Synth.} & \textbf{Fine-tune} & \textbf{Notes} \\
\midrule
\code{img\_height}               & 64                 & 64                 & half the melody staff height \\ \tfgrowsep
\code{rnn\_hidden}               & 192                & 192                & per direction; smaller than melody CRNN (256) \\ \tfgrowsep
\code{rnn\_layers}               & 2                  & 2                  & BiLSTM depth \\ \tfgrowsep
\code{dropout}                   & 0.2                & 0.2                & between LSTM layers \\ \tfgrowsep
\code{cnn\_dropout}              & 0.15               & 0.15               & inter-block CNN dropout \\ \tfgrowsep
\code{batch\_size}               & 32                 & 32                 & \\ \tfgrowsep
\code{lr}                        & $1{\times}10^{-3}$ & $2{\times}10^{-4}$ & OneCycleLR peak; reduced for fine-tuning \\ \tfgrowsep
\code{epochs}                    & 30                 & 20                 & maximum; early stopping active in both phases \\ \tfgrowsep
\code{early\_stopping\_patience} & 8                  & 8                  & epochs without CER improvement \\ \tfgrowsep
\code{synth\_weight}             & N/A                & 0.4                & weight of synthetic strips in the fine-tune sampler \\
\bottomrule
\end{tabularx}
\caption{Chord CRNN hyperparameters for the two training phases.
The fine-tuning learning rate is five times lower to avoid
displacing representations learned on the larger synthetic corpus.
\code{synth\_weight} controls the relative frequency of synthetic
versus real strips in the fine-tuning batch; real strips always
carry weight~1.0.}
\label{tab:chord-hyperparams}
\end{table}

\subsection{Fine-tuning on hand-labelled strips}

After synthetic training, the model is fine-tuned on a small set of
real chord strips extracted from Real Book pages and manually
labelled.
Extraction reuses the staff-detection module of the inference
pipeline: chord strips are cropped from the region directly above each
detected staff, and the synthetic chord model is run on each strip to
provide a first-pass label that the human reviewer corrects rather
than typing from scratch.
Labels are stored in a JSON Lines file with a \code{status} field
(\code{pending} / \code{done} / \code{skip}) so that the labeller can
skip strips that contain no chords or are too degraded to read; the
hand-labelling UI is served at \code{/labeler} on the FastAPI server
(\Cref{sec:impl-api}).

Fine-tuning mixes real and synthetic strips with a weighted sampler:
real strips carry weight $1.0$, synthetic strips weight $0.4$ (the
default \code{--synth-weight}).  The synthetic admixture prevents
catastrophic forgetting of rare chord qualities under-represented in
the small real set.
Labels go through the \term{canonicalisation} step defined in
\Cref{sec:design-twostream} before the vocabulary filter:
\code{-7b5}/\code{m7b5}/\code{min7b5} are rewritten to the
canonical \code{ø} token; \code{Cm7} becomes
\code{C-7}; \code{sus4}/\code{sus2} become \code{sus}; and any chord
token containing an out-of-vocabulary character is dropped (with a
warning) rather than silently corrupted by mid-token character
removal.

% ─────────────────────────────────────────────────────────────────────────────
\section{Inference Pipeline Implementation}
\label{sec:impl-inference}

The five-stage data flow was established in \Cref{sec:design-pipeline}
and illustrated in \Cref{fig:pipeline-design}; the present section
describes how each stage is realised in code.
The pipeline is a single function, \code{run\_pipeline}, that accepts
the raw bytes of a PDF or image plus the original filename and returns
a structured JSON dictionary.
It is invoked by the FastAPI endpoint of \Cref{sec:impl-api}; no
separate CLI subcommand is exposed for single-page inference because
the typical use-case is uploading a page through the web UI.

\subsection{Preprocessing}

PDFs are rasterised with PyMuPDF at the DPI specified by
\code{OMR\_PDF\_DPI} (default \num{300}, clamped to $[72, 600]$).
The current implementation processes only page~0 of multi-page PDFs;
multi-page batching is left to future work.
Raster images are loaded directly via OpenCV.
Binarisation applies Contrast Limited Adaptive Histogram Equalisation
(CLAHE) followed by Otsu thresholding; CLAHE first normalises the
local contrast so that the global Otsu threshold becomes robust to
uneven scanner illumination.
Deskewing tests a range of rotation angles and selects the one that
maximises the variance of the horizontal projection profile
(sharper row peaks correspond to better-aligned staves).

\subsection{Staff detection}

Staff detection is implemented as a classical morphological pipeline,
in line with the design argument of \Cref{sec:design-staffdetect}.
A wide horizontal opening kernel extracts staff-line candidates, row
centroids are computed and merged across closely-spaced rows to handle
thick lines, and the centroids are clustered into groups of five
consecutive nearly-equispaced lines.
For each cluster, the system emits a \code{System} record containing
the cropped staff image, the strip of pixels above it (vertical span
computed as the gap to the previous staff), and the bounding box on
the page.

\subsection{Rejection gates}

Between staff detection and the CRNN, every candidate strip is passed
through the hybrid pre/post-CRNN gate of
\Cref{sec:design-staffdetect}.
The pre-CRNN gate evaluates five geometric and OCR-based signals:
local re-detection of five staff lines
(\code{geometry\_no\_staff\_lines}),
per-line horizontal coverage (\code{geometry\_line\_span}),
coefficient of variation of inter-line gaps
(\code{geometry\_spacing\_cov}),
inter-line ink density (\code{geometry\_interline\_ink}),
and EasyOCR text-box area fraction
(\code{ocr\_text\_density}).
The OCR signal is provided by a lazily-loaded EasyOCR reader cached at
module scope; the reader is configured for CPU because the surrounding
gate runs once per detected strip and the CPU latency is bounded.

Strips that survive the pre-CRNN gate (and borderline strips
re-admitted via the CTC-confidence override) are then evaluated by
the post-CRNN gate, which thresholds the mean log-probability of the
CRNN's argmax frames against the calibrated
\code{min\_mean\_logprob}.
Thresholds live in \code{models/staff\_reject/thresholds.json} and
are produced by the \code{calibrate-reject} CLI subcommand
(\Cref{sec:impl-cli}); the file is checkpoint-specific and must be
regenerated whenever the music CRNN is retrained.

\subsection{Music recognition}

Each \code{System} is fed through the staff-aware normalisation
(resize to \SI{128}{px} height, centre the staff vertically inside the
frame, per-image zero-mean and unit-variance), then batched and passed
to the CRNN.
The model is loaded once per process and cached in module scope; the
checkpoint path is resolved from \code{models/latest/best\_model.pt}
or from the \code{OMR\_CHECKPOINT} environment variable.
Images wider than \num{2048}~px are clamped at load time to avoid
out-of-memory failures on very wide single staves.

Decoding is greedy by default: argmax over the vocabulary dimension at
each time step, collapse consecutive duplicates, remove blank.
Beam search is opt-in via the \code{OMR\_BEAM\_WIDTH=N} environment
variable or the \code{--beam-width N} CLI flag.
The performance comparison between the two modes is reported in
\Cref{ch:results}; the present chapter only commits to greedy as the
runtime default and to the implementation cost of beam search, which
is examined in \Cref{sec:impl-performance}.

\subsection{Chord recognition}

The chord strip from each system is trimmed at its left edge to
remove binder-hole shadows (up to \SI{10}{\percent} of strip width,
controlled by a column-wise ink-density threshold of $\geq 50\%$),
resized to the chord network's training height of \SI{64}{px},
batched, and decoded.
\code{<unk>} tokens emitted when the network encounters clutter the
synthetic training set never saw are treated as word separators
rather than as character errors.
The decoded character string is then passed to a jazz-chord grammar
filter (\code{chord\_postprocess.clean\_chord\_line}) that validates
against \code{ROOT [ACC] [QUALITY] [EXTEN] [ALT]* [SLASH]} and drops
fragments that do not match (page numbers, tempo markings, stray
lyrics).
Single-character predictions consisting of a root letter only
(no quality, no extension) are also dropped as likely false
positives.

\subsection{Grammar fixer}

The grammar fixer is implemented as a stateful walker over the raw
token list, realising the design contract of \Cref{sec:design-grammar}.
The walker tracks the accumulated beat count against the active time
signature, the active key signature, and whether a pitch token has
been emitted that is still awaiting its octave and duration.
The concrete state-transition rules applied per token are:

\begin{itemize}
    \item A pitch token starts a new note; the walker waits for an
          octave token and then a duration token before committing
          the note to the output stream.
    \item A duration token without a preceding pitch is rejected.
    \item A duration appearing immediately after another duration
          ends the prior note prematurely.
    \item An accidental appearing before the duration is reordered to
          appear after it.
    \item Octaves outside $[3, 7]$ are clamped to the range.
    \item Clefs other than \code{clef:G2} are rewritten.
    \item Time signatures outside the eight-element allowed set are
          dropped, and the prior measure's signature is propagated.
    \item \code{tied:stop} tokens are kept only if a matching
          \code{tied:start} appeared earlier in the same tie window.
\end{itemize}

The allowed time-signature set
(\code{\{4/4, 3/4, 2/4, 2/2, 6/8, 6/4, 5/4, 12/8\}}) is stored in
both \code{src/CRNN\_CTC/dataset.py} and
\code{src/omr\_pipeline/grammar\_fix.py}; the two definitions must
stay synchronised (Hard~Constraint~5 of the contributor guide).
The walker emits a normalised LMX token list that satisfies the
grammar of \Cref{sec:lmx-format}.

% ─────────────────────────────────────────────────────────────────────────────
\section{Command-Line Interface}
\label{sec:impl-cli}

Every offline operation goes through \code{src/cli.py}, a single
\code{argparse} dispatcher with the twelve subcommands listed in
\Cref{tab:cli-subcommands}.

\begin{table}[htbp]
\centering
\small
\setlength{\tabcolsep}{5pt}
\renewcommand{\arraystretch}{1.35}
\begin{tabularx}{\textwidth}{@{} l L @{}}
\toprule
\textbf{Subcommand} & \textbf{Purpose} \\
\midrule
\code{render}                 & Render PrIMuS \code{.semantic} files to LilyJAZZ PNGs. \\ \tfgrowsep
\code{convert}                & Convert PrIMuS \code{.semantic} files to LMX labels. \\ \tfgrowsep
\code{augment}                & Apply scan-simulation augmentation to clean PNGs. \\ \tfgrowsep
\code{vocab}                  & Build the LMX vocabulary file from a corpus. \\ \tfgrowsep
\code{pipeline}               & Run \code{render} + \code{convert} + twins + \code{augment} + \code{vocab} (\code{--force-all} rebuilds all). \\ \tfgrowsep
\code{train}                  & Train the CRNN on a prepared dataset. \\ \tfgrowsep
\code{evaluate}               & Evaluate a checkpoint and report SER. \\ \tfgrowsep
\code{evaluate-ab}            & Compare beam widths on the same checkpoint. \\ \tfgrowsep
\code{api}                    & Launch the FastAPI web service. \\ \tfgrowsep
\code{pipeline-train}         & Run the full data pipeline and then immediately start training. \\ \tfgrowsep
\code{harvest-reject-fixtures} & Cut staff-strip fixtures from real PDFs for reject-gate calibration. \\ \tfgrowsep
\code{calibrate-reject}       & Sweep reject-gate thresholds on labelled fixtures and emit \code{thresholds.json}. \\
\bottomrule
\end{tabularx}
\caption{CLI subcommands exposed by \code{src/cli.py}.  Every
subcommand prints a self-contained help text when invoked with
\code{--help}.}
\label{tab:cli-subcommands}
\end{table}

CLI flags map one-to-one onto fields of the \code{Config} dataclass for
training and evaluation subcommands, so a flag added to \code{Config}
becomes a flag of the corresponding subcommand without further
plumbing.

% ─────────────────────────────────────────────────────────────────────────────
\section{Web Service}
\label{sec:impl-api}

The system exposes a FastAPI~+~Uvicorn web service that wraps
\code{run\_pipeline} as a REST endpoint.
The principal route, \code{POST /api/omr/lead-sheet}, accepts a
multipart file upload (PDF or image) and returns the JSON dictionary
produced by the pipeline; a \code{GET /health} route reports
liveness and the resolved checkpoint path, suitable as a container
healthcheck probe.
Two additional UIs support hand-labelling of training data:
\code{/labeler} for chord strips and \code{/music-labeler} for staff
strips, together with their supporting \code{/api/labeler/*} and
\code{/api/music-labeler/*} routes used during the fine-tuning
workflows of \Cref{sec:impl-chord}.

The CRNN models are loaded lazily on the first inference request and
cached for the lifetime of the process; the server runs on a single
Uvicorn worker by default to keep GPU contention out of the deployment
contract.
A minimal HTML upload UI is served at the root path (\code{GET /});
it displays the rendered page with detected staff bounding boxes
overlaid, the recognised LMX token sequence per segment, and the
extracted chord list.
The UI is a debugging aid that doubles as a demonstration interface;
it is not the principal deliverable.
\todo{Run \texttt{notebooks/02\_evaluate\_model.ipynb} §9.2 to get
median page latency; replace this line with the result.}
The full response schema is documented in
\code{docs/api.md}; the JSON top level carries \code{pages[]},
\code{meta\{\}}, and a top-level \code{error} field that is non-null
only on pipeline failure.

% ─────────────────────────────────────────────────────────────────────────────
\section{Computational Performance and Speed Optimisations}
\label{sec:impl-performance}

This section covers the computational-complexity reasoning behind the
design of \Cref{ch:design} and the concrete engineering choices that
follow from it: which operations are cheap by construction, which
required active work to keep within budget, and how the
throughput/latency/accuracy trade-offs are exposed to the operator.

\subsection{Algorithmic-complexity choices}
\label{sec:impl-perf-algorithmic}

\paragraph{CTC decoding: greedy vs beam search.}
Greedy CTC decoding has complexity $\mathcal{O}(T\cdot V)$ per
sample, where $T$ is the number of output time steps and $V$ is the
vocabulary size: it consists of one argmax per time step followed by
a linear pass to collapse duplicates and remove blanks.
Beam search of width $k$ has complexity $\mathcal{O}(T\cdot V\cdot k)$
with an additional logarithmic factor for the priority queue.
On the monophonic Real Book corpus, beam search of width $5$ yields
$\leq 1$ edit per $1000$ SER improvement at approximately $6\times$
the decode cost; greedy is therefore the runtime default for both
the CLI \code{evaluate} subcommand and the API.
Beam search is opt-in for users who want to spend the latency to
chase the small additional accuracy on dense passages.

\paragraph{Width compression and the time axis.}
The $4\times$ width compression in the CNN backbone
(\Cref{tab:resnet-strides}) directly controls the cost of every
downstream operation: the BiLSTM, the linear projection, and the CTC
forward--backward all scale linearly in $T \approx W/4$.
A smaller compression ratio would multiply the recurrent and projection
costs by the same factor; a larger ratio risks producing more output
tokens than time steps, which CTC's \code{zero\_infinity} option
silently masks but which corrupts the gradient estimate.
The chosen $4\times$ ratio is the smallest power-of-two ratio at which
every in-scope sample satisfies $T \geq L$ for label length $L$, while
keeping the time axis rich enough to disambiguate adjacent symbols on
densely-engraved staves.

\paragraph{ResNet18 vs VGG (parameter count and FLOPs).}
The ResNet18 backbone has approximately
\num{11.2}~million parameters and a
forward-pass cost of roughly \num{1.8}~GFLOP at the project's
input shape (both figures are estimates for the standard ResNet18 at
comparable input resolution; exact values vary with the asymmetric
stride modifications described above).
The legacy VGG-style stack at the same number of channels carries
more parameters and runs slower without measurable accuracy gain on
this corpus; that comparison is the empirical reason ResNet18 became
the default backbone, as discussed in \Cref{sec:design-model}.
ResNet18's residual connections also help the gradient propagate
deep into the network during the early epochs of CTC training, where
the model is still learning the rough alignment between input frames
and output tokens.

\paragraph{BiLSTM vs Transformer at this scale.}
The two-layer BiLSTM with hidden size \num{256} per direction has on
the order of \num{2.6}~M parameters and runs in
$\mathcal{O}(T\cdot d^2)$ per layer, with $d = 512$ as the combined
hidden width.
A Transformer encoder layer over the same time axis runs in
$\mathcal{O}(T^2\cdot d)$ for the self-attention, plus
$\mathcal{O}(T\cdot d^2)$ for the feed-forward block.
For typical staff widths in this corpus, $T \approx 200$--$500$, so
self-attention is not asymptotically prohibitive, but a Transformer
also requires substantially more parameters and training data to
reach competitive accuracy on a monophonic task
(\Cref{sec:design-model-alts}).
The BiLSTM was preferred because it converges in fewer epochs on a
synthetic corpus of this size and fits the project's compute envelope.

\paragraph{Monophonic-only output space.}
The single most consequential complexity decision is the
monophonic-only contract (\Cref{sec:design-domain}): a single melody
line per staff turns the recognition problem into a strictly
left-to-right CTC alignment, with no need for tree decoding or
position-aware attention over polyphonic voices.
The output vocabulary collapses from the thousands of tokens needed by
polyphonic systems~\cite{riosVilaSheetMusicTransformer2024} to
$\approx 80$ symbols on a full corpus rebuild (empirical union of
\code{.lmx} tokens plus guaranteed pitch/octave ranges), and the CTC
forward--backward shrinks proportionally.
This is the lever that makes a single-GPU Bachelor's thesis budget
viable on a task that otherwise demands an order of magnitude more
compute.

\subsection{Engineering optimisations}

\paragraph{Process-based parallelism for the data pipeline.}
The five offline stages (\code{render}, \code{convert}, twins,
\code{augment}, \code{vocab}) are CPU-bound: LilyPond invocation, image augmentation
under \code{albumentations}, and string parsing dominate their cost,
none of which release the Python Global Interpreter Lock for long
enough to make threading useful.
The implementation uses \code{multiprocessing.Pool} with a default of
\num{8} worker processes (configurable via \code{--workers}); on the
target machine this saturates the available CPU cores during dataset
generation.
A thread-based alternative was rejected by direct measurement during
development: the threaded variant was slower because of GIL
contention on the per-sample LilyPond subprocess wrapper.

\paragraph{Mixed-precision training (AMP).}
The training loop wraps the forward pass in
\code{torch.amp.autocast(device\_type=}\allowbreak\code{"cuda")} and
scales the loss with \code{torch.amp.GradScaler}.
Numerically delicate operations (the CTC forward--backward,
log-softmax) are kept in single precision, while the convolutional
and LSTM bodies run in \code{float16}.
The principal benefit on the target hardware is VRAM rather than
throughput: AMP keeps peak memory below
\SI{8.5}{\giga\byte} at \code{batch\_size=16}, leaving headroom for
the data-loader workers and the validation pass.

\paragraph{Batch size against the VRAM ceiling.}
The default batch size of \num{16} was selected as the largest value
that fits within \SI{12}{\giga\byte} of VRAM at the target image
height of \SI{128}{px}, with AMP enabled and \code{num\_workers=10}
on the data-loader.
Batches of \num{24} are feasible on the target hardware at the cost
of slightly tighter memory headroom; batch \num{32} requires either
disabling AMP-safe activations or reducing image height.
Gradient accumulation was not implemented because the CTC gradient is
well-behaved at the current batch size; the LSTM gradient is the
primary source of instability and is controlled by the global-norm
clip of \num{5.0} rather than by a larger effective batch.
A larger effective batch would in principle improve gradient noise but
would also lengthen the OneCycleLR cycle and slow time-to-convergence.

\paragraph{DataLoader configuration.}
The training and validation loaders use \code{num\_workers=10},
\code{pin\_memory=True}, and \code{persistent\_workers=True}.
The first parameter sets the number of CPU processes feeding the
GPU; ten was chosen empirically as the value at which the GPU is no
longer waiting for I/O on the target machine.
\code{pin\_memory} accelerates host-to-device transfers by allocating
pinned memory buffers, and \code{persistent\_workers} keeps the workers alive between epochs,
avoiding the cost of respawning ten processes at each epoch boundary;
the improvement was noticeable during development but was not
formally profiled.

\paragraph{Fused AdamW.}
On CUDA, the optimiser is constructed with \code{fused=True}, which
collapses the parameter-update kernel into a single CUDA launch per
parameter group rather than one launch per tensor.  The saving is
small in absolute terms (a few hundred milliseconds per epoch) but
free, since the fused implementation is identical numerically.

\paragraph{Inference model caching.}
Both CRNN checkpoints (music and chord) are loaded into memory once
per process and reused across calls.
The first request to the API pays the full checkpoint-load cost,
later requests do not.
For the FastAPI deployment this matters because the typical
deployment runs a single Uvicorn worker process per GPU; cold-start
latency is paid once at process start, not per request.

\subsection{Inference latency}

\paragraph{End-to-end per-page wall time.}
On the target RTX~3060, the end-to-end inference cost of a typical
\num{8}-staff Real Book page is approximately
\todo{Run \texttt{notebooks/02\_evaluate\_model.ipynb} §9.2
and paste the median from the summary cell here.}.

\paragraph{Cost breakdown across stages.}
The morphological staff detection runs in milliseconds on CPU and is
not the bottleneck for the typical input.
The pre-CRNN rejection gate dominates the latency budget for the
non-music regions: EasyOCR's text detector is the most expensive
component because it runs the full CRAFT detection model on each
candidate strip.
Once a strip is accepted, the music CRNN forward pass and greedy CTC
decoding together cost roughly
\todo{Run \texttt{notebooks/02\_evaluate\_model.ipynb} §9.1
and paste the median per-strip latency here.}.
The chord CRNN runs on a strip of half the staff height and a small
fraction of the strip width, so its forward-pass cost is correspondingly
smaller; its total throughput is dominated by the per-batch overhead of
launching CUDA kernels.
The grammar fixer is a pure-Python state machine over the predicted
token sequence; its cost is negligible.

\subsection{Throughput, latency, and accuracy trade-offs}

The system exposes three knobs that move along the throughput /
latency / accuracy frontier.
\Cref{tab:perf-tradeoffs} summarises them.

\begin{table}[ht]
\centering
\small
\setlength{\tabcolsep}{5pt}
\renewcommand{\arraystretch}{1.30}
\begin{tabularx}{\textwidth}{@{} l L L L @{}}
\toprule
\textbf{Knob} & \textbf{Throughput effect} & \textbf{Latency effect} & \textbf{Accuracy effect} \\
\midrule
Beam width $k$ (\code{OMR\_BEAM\_WIDTH}) & $\sim 1/k$ at fixed throughput per stripe. & $\sim k\times$ per strip. & $\leq 1$ edit per $1000$ SER improvement at $k=5$ on this corpus. \\ \tfgrowsep
PDF DPI (\code{OMR\_PDF\_DPI}) & Linear: image area scales as DPI$^2$. & Linear in image area. & Higher DPI improves CRNN accuracy up to roughly \SI{300}{DPI}; beyond that the gain saturates and pre-processing dominates. \\ \tfgrowsep
Pre-CRNN reject gate & Improves throughput by skipping non-music strips' CRNN call. & Reduces the average latency per page. & Trade-off between false-positive rejection of valid staves and false-negative acceptance of text; thresholds are calibrated on labelled fixtures. \\
\bottomrule
\end{tabularx}
\caption{Operator-tunable trade-offs in the inference pipeline.  Exact
multipliers are stated where measured in this thesis; the SER /
latency numbers themselves are reported in \Cref{ch:results}.}
\label{tab:perf-tradeoffs}
\end{table}

The default configuration of the system commits to greedy decoding,
\SI{300}{DPI} rasterisation, and the calibrated rejection thresholds
in \code{models/staff\_reject/thresholds.json}: this is the operating
point that minimises latency while remaining within the accuracy
targets stated in \Cref{sec:design-domain}.
Other operating points are reachable by setting the corresponding
environment variable without retraining or recompiling.

% ─────────────────────────────────────────────────────────────────────────────
\section{Reproducibility}
\label{sec:impl-reproducibility}

Three mechanisms together guarantee that any experiment reported in
\Cref{ch:results} can be reproduced exactly.

\paragraph{Deterministic dataset splits.}
The \code{make\_splits} function partitions samples by hashing the
sample identifier and bucketing the hash modulo a fixed denominator.
The same sample always falls into the same split across runs and
across machines, so any difference between two evaluation runs comes
from training, not from the split.

\paragraph{Embedded configuration.}
Every checkpoint contains a JSON-serialised copy of the \code{Config}
dataclass that produced it.
At evaluation or inference time, the embedded config is loaded before
the model weights, so that the data filters, image dimensions, and
hyperparameters are exactly those used during training.
This is also why the \code{Config} dataclass is a hard contract:
adding, renaming, or removing a field breaks backward compatibility
with existing checkpoints (Hard~Constraint~3 of the project's
contributor guide).

\paragraph{Pinned dependencies.}
Poetry's \code{poetry.lock} file pins every dependency to an exact
version, including transitive ones, and \code{pyproject.toml} fixes
the Python version at \num{3.14}.
The project environment is reproducible across machines with a single
\code{poetry install}.
The lock file is committed to version control alongside the source.
The corpus was rendered with \textbf{LilyPond~2.26.0} (Guile~3.0
runtime).
A future LilyPond release could shift rendered pixels even for
identical source files, so the version is recorded here as part of
the reproducibility chain.
The LilyJAZZ font and stylesheet are installed from the
OpenLilyPondFonts repository into the LilyPond \code{share} directory
during setup; the installation procedure, including the API patch
required for LilyPond~$\geq$~2.25, is documented in
\code{docs/data\_pipeline.md}.

\paragraph{Random-seed strategy.}
The training script calls \code{seed\_everything(cfg.seed)} at
startup before any data is loaded.
This function seeds all four sources of randomness in one place:
\code{random.seed}, \code{numpy.random.seed}, \code{torch.manual\_seed},
and \code{torch.cuda.manual\_seed\_all}.
It also sets \code{torch.backends.cudnn.deterministic=True} and
\code{torch.backends.cudnn.benchmark=False}: the first forces
deterministic convolution paths; the second disables CuDNN's
auto-tuner, which would otherwise select different kernels from run to
run based on input shapes.
The seed value defaults to \num{42} and is a field of \code{Config},
so it is serialised inside every checkpoint alongside the rest of the
training parameters.

Dataset splitting uses a separate \code{torch.Generator} seeded from
\code{Config.seed}, keeping the train/val/test partition deterministic
independently of any global state.
The offline augmentation stage accepts a \code{--seed} CLI flag that
seeds the per-sample randomness in \code{albumentations}; passing the
same seed and inputs reproduces identical augmented images on a second
run of the pipeline.
One caveat applies: a small number of CUDA reduction operations
involved in the \code{zero\_infinity} CTC path do not have
deterministic implementations on all GPU architectures, so
bitwise-identical outputs across different GPU generations cannot be
guaranteed even with the above settings.

\paragraph{Checkpoint contents.}
Each saved checkpoint is a PyTorch dictionary with keys
\code{model\_state\_dict}, \code{optimizer\_state\_dict},
\code{epoch}, \code{val\_ser}, and \code{config}.  The optimiser
state is preserved so that \code{--resume} restores the AdamW
moments and the OneCycleLR schedule from the last completed epoch.
The corresponding \code{training\_log.csv} records per-epoch
training loss, validation loss, and validation SER; together with
the checkpoint these files reconstruct the training trajectory in
full.
